\documentclass{article}
\usepackage{amsmath, amssymb, amsthm}
\usepackage{hyperref}

\title{Welcome to texview}
\author{A browser-only LaTeX viewer}
\date{}

\newtheorem{theorem}{Theorem}
\newtheorem{lemma}{Lemma}
\newtheorem{definition}{Definition}

\begin{document}
\maketitle

A browser-only viewer for lightweight LaTeX. Aimed at math notes, problem
sets and derivations rather than full papers --- there's no bibliography
engine, no real float placement and no macro expansion, but math,
sectioning, lists, theorem-like environments, links, \texttt{verbatim},
basic \texttt{tabular} and \texttt{\textbackslash includegraphics} all
render live as you type.

\textit{Tip: press \texttt{?} for the keyboard cheat sheet, or
\texttt{Alt+G} to reload this guide.}

\section{What you can do}

\begin{itemize}
  \item \textbf{Side-by-side preview} with synchronized scrolling
  \item \textbf{Selection mirroring}: highlight text in either pane and
        the other follows
  \item \textbf{Four view modes:} split, stacked, source only, preview only
  \item \textbf{Math, lists, tables, theorems, code, figures}
  \item \textbf{Drag-and-drop} or paste TeX source to load a file
  \item \textbf{Image attachments} with figure-stub insertion
  \item \textbf{Export} as HTML, print to PDF, or download the \texttt{.tex}
  \item \textbf{Persists} across reloads (your last session is restored)
  \item \textbf{Light / dark theme} with auto-detection
  \item \textbf{Mobile-friendly} layout that adapts to narrow screens
\end{itemize}

\section{Math}

Inline math uses \texttt{\$\dots\$} or \texttt{\textbackslash(\dots\textbackslash)}:
the area of a circle is $A = \pi r^2$, Euler's identity is
$e^{i\pi} + 1 = 0$, and the golden ratio satisfies
$\varphi^2 = \varphi + 1$.

\subsection{Display}

Display math uses \texttt{\$\$\dots\$\$},
\texttt{\textbackslash[\dots\textbackslash]}, or any \texttt{amsmath}
environment KaTeX supports:

\[
  \frac{\partial}{\partial t}\Psi(\mathbf{r}, t)
  = \frac{1}{i\hbar}\,\hat{H}\,\Psi(\mathbf{r}, t)
\]

\subsection{Aligned systems}

\begin{align}
  \nabla \cdot \mathbf{E}  &= \frac{\rho}{\varepsilon_0} \\
  \nabla \cdot \mathbf{B}  &= 0 \\
  \nabla \times \mathbf{E} &= -\frac{\partial \mathbf{B}}{\partial t} \\
  \nabla \times \mathbf{B} &= \mu_0 \mathbf{J}
                            + \mu_0\varepsilon_0\frac{\partial \mathbf{E}}{\partial t}
\end{align}

\subsection{Cases}

\[
  \mathrm{sign}(x) = \begin{cases}
    +1 & x > 0 \\
    \phantom{+}0 & x = 0 \\
    -1 & x < 0
  \end{cases}
\]

\subsection{Matrices}

\[
  R(\theta) =
  \begin{pmatrix}
    \cos\theta & -\sin\theta \\
    \sin\theta & \phantom{-}\cos\theta
  \end{pmatrix},
  \qquad
  R(\theta)\,R(\varphi) = R(\theta + \varphi).
\]

\subsection{Sums, products, integrals}

\[
  \sum_{n=1}^{\infty} \frac{1}{n^2} = \frac{\pi^2}{6},
  \qquad
  \prod_{p \text{ prime}} \frac{1}{1 - p^{-s}} = \zeta(s),
  \qquad
  \int_{-\infty}^{\infty} e^{-x^2}\,dx = \sqrt{\pi}.
\]

\section{Sectioning}

This is a paragraph. The \texttt{\textbackslash section} command takes
\textbf{a numbered heading}; \texttt{\textbackslash section*} gives an
\textit{unnumbered} one.

\subsection{Subsections}

Section numbers nest. Cross-references resolve too --- e.g.\ this lives
in section \ref{sec:numbers}.

\subsection{Numbers and refs}\label{sec:numbers}

You can attach \texttt{\textbackslash label\{\dots\}} to a heading and
later cite it with \texttt{\textbackslash ref\{\dots\}} or
\texttt{\textbackslash autoref\{\dots\}}; \texttt{\textbackslash eqref}
renders the parenthesised equation number.

\subsubsection{Even deeper}

Down to \texttt{\textbackslash subsubsection};
\texttt{\textbackslash paragraph} and \texttt{\textbackslash subparagraph}
also work, but render as small inline-style headings.

\section{Lists}

Unordered:
\begin{itemize}
  \item apples
  \item pears
  \item citrus
\end{itemize}

Ordered:
\begin{enumerate}
  \item wake up
  \item make coffee
  \item read \texttt{.tex} files
\end{enumerate}

Description:
\begin{description}
  \item[\texttt{KaTeX}] fast math typesetting in the browser
  \item[\texttt{CodeMirror}] the source editor
  \item[\texttt{DOMPurify}] HTML sanitiser
\end{description}

\section{Theorem-like environments}

\begin{theorem}[Pythagoras]
For any right triangle with legs $a$, $b$ and hypotenuse $c$,
\[ a^2 + b^2 = c^2. \]
\end{theorem}

\begin{lemma}
The sum of the first $n$ positive integers is
$\tfrac{n(n+1)}{2}$.
\end{lemma}

\begin{proof}
Pair the $k$-th term with the $(n+1-k)$-th; each pair sums to $n+1$,
and there are $n/2$ such pairs (or $(n-1)/2$ pairs and a middle term
when $n$ is odd; either way the total is $n(n+1)/2$).
\end{proof}

\begin{definition}
A function $f \colon \mathbb{R} \to \mathbb{R}$ is
\emph{continuous at $x_0$} if for every $\varepsilon > 0$ there is a
$\delta > 0$ such that $|x - x_0| < \delta$ implies
$|f(x) - f(x_0)| < \varepsilon$.
\end{definition}

\section{Code}

Inline code with \verb|\verb|, like
\verb|cm.setOption('lineNumbers', false)| or
\verb!sed -e 's/foo/bar/g' file.txt!.

\begin{verbatim}
function fibonacci(n) {
  if (n < 2) return n;
  let [a, b] = [0, 1];
  for (let i = 2; i <= n; i++) [a, b] = [b, a + b];
  return b;
}
\end{verbatim}

\section{Tables}

A simple \texttt{tabular}:

\begin{center}
\begin{tabular}{l c r}
  \hline
  Feature       & Library      & Notes                     \\
  \hline
  Math          & KaTeX        & fast, no MathJax          \\
  Editor        & CodeMirror 5 & \texttt{stex} mode        \\
  Sanitisation  & DOMPurify    & always-on safety net      \\
  \hline
\end{tabular}
\end{center}

\section{Inline formatting}

Mix \textbf{bold}, \textit{italic}, \emph{emphasised}, \texttt{monospace},
\textsf{sans-serif}, \textsc{small caps}, \underline{underlined} and
\sout{struck-through} text. Subscripts and superscripts also work in
running text: H\textsubscript{2}O, x\textsuperscript{2}. The TeX logo:
\TeX\ and \LaTeX.

\section{Quotes}

\begin{quote}
``Any sufficiently advanced technology is indistinguishable from magic.''
--- Arthur C.\ Clarke
\end{quote}

The renderer turns \verb|---| into an em-dash, \verb|--| into an
en-dash, and \verb|``...''| into curly quotes.

\section{Links}

External: \href{https://katex.org/}{KaTeX} powers math rendering;
\href{https://codemirror.net/5/}{CodeMirror~5} powers the editor.
Bare URLs use \texttt{\textbackslash url}:
\url{https://github.com/utksi/texview}.

\section{Figures and attachments}

Drop an image file anywhere on the page (or paste a screenshot from the
clipboard), then click the paperclip in the toolbar and pick
\textbf{Insert reference} --- a \texttt{figure} block is inserted at
the cursor:

\begin{verbatim}
\begin{figure}
  \centering
  \includegraphics[width=0.6\textwidth]{./diagram.png}
  \caption{Sketch of the setup}
\end{figure}
\end{verbatim}

On download, the \texttt{.tex} and every attachment are saved together;
keep them in the same folder so the \texttt{./name.ext} references
resolve.

\section{What's not supported}

Bibliographies (\texttt{\textbackslash cite},
\texttt{\textbackslash bibliography}), custom macros
(\texttt{\textbackslash newcommand} is stripped, not expanded), exact
float placement, page-level layout, fine microtypography, and most
package-specific commands. If you need any of those, a real LaTeX engine
(local TeX Live, Overleaf, or a WASM build like SwiftLaTeX) is the right
tool.

\section*{Try it}

Load your own file (open icon in the toolbar, or drop one anywhere),
or just start typing in the left pane. Switch view modes with
\texttt{Alt+1} through \texttt{Alt+4}. Drag the divider to resize the
panes; double-click it to reset to 50/50.

\end{document}
